\chapter{Análisis del sistema}
\label{anexo2:cap:analisis}
\textbf{TODO}

\section{Modelo del dominio}
\label{anexo2:sec:modelo-dominio}
El modelo del dominio es una representación conceptual de los objetos y las relaciones entre ellos
en el sistema. Representa, de manera muy general y sin entrar en detalles de implementación, tanto
el comportamiento como los atributos de los objetos que forman parte del sistema
\parencite{DomainModel}.

De manera muy general, el funcionamiento del sistema se puede resumir en: los usuarios del sistema
se registran y pueden modificar sus datos personales. Estos usuarios pueden registrar su forma física,
así como modificarlas y consultarlas. Además, el sistema ofrece una serie de actividades que los usuarios
pueden realizar, y que una vez realizadas quedan registradas en el sistema. Los usuarios pueden
consultar las actividades que han realizado, así como las estadísticas de su forma física.

\diagrama{
    nombre = {Modelo del dominio},
    archivo = {anexo2/diagramas/ModeloDominio.pdf},
    etiqueta = {anexo2:fig:modelo-dominio},
    forzarPosicion = {sí pls},
}
